% TOP_PROBLEM: {"id": 30002082, "title": "Uniqueness of Tangent Flows in Mean Curvature Flow", "category_id": 9, "category": {"id": 9, "name": "pde", "display_name": "Partial Differential Equations", "description": "PDEs and their applications in physics and geometry.", "slug": "pde", "order_index": 9, "created_at": "2026-07-31T15:26:25.671Z"}, "status": "partially_solved"}
% FIRST_POSTED: 2026-10-06
% ATTEMPT_STATUS: unresolved
% OPENAI_CLAIM: {"schema_version":1,"source_repo":"https://github.com/openai/math","source_commit":"adc7f1241b42e322a6451854ab7e4b4c146bf78a","release_date":"2026-10-06","match":"partial","resolution":"partial","independently_reviewed":false,"second_pass":"n/a","families":[{"family":"355","title":"Unique tangent flows at the first surface singularity","subject":"Differential geometry","manuscripts":[{"title":"Uniqueness of tangent flows at the first singular time of embedded surface mean-curvature flow","pdf_path":"preprints/Tangent-flow-uniqueness-2026-09-24/paper.pdf","readme_path":"preprints/Tangent-flow-uniqueness-2026-09-24/README.md","date":"2026-09-24","theorem_ref":"build/sections/introduction.tex:46, Theorem thm:main"}],"lean":null,"reasoning_trace":null}]}
\documentclass[11pt]{article}
\usepackage[margin=25mm]{geometry}
\usepackage{fontspec}
\setmainfont{Latin Modern Roman}
\usepackage{amsmath,amssymb,mathtools}
\usepackage{unicode-math}
\setmathfont{Latin Modern Math}
\usepackage{enumitem}
\usepackage{xurl,hyperref,bookmark}
\hypersetup{colorlinks=true,urlcolor=blue}
\setcounter{secnumdepth}{0}
\setlength{\parindent}{0pt}
\setlength{\parskip}{5pt}
\setlength{\emergencystretch}{3em}
\begin{document}
\section*{Uniqueness of Tangent Flows in Mean Curvature Flow}

\subsection{OpenAI's claim}
\textit{Uniqueness of surface tangent flow}: At the first singular time of a smooth closed embedded surface flow in $\mathbb R^3$, the paper claims a unique multiplicity-one self-shrinking tangent flow at each fixed singular point. Its position and axes are unique in the original coordinates.

\subsection{Outline of the argument}
\textit{Uniqueness of surface tangent flow}: The stationary analysis handles conical and cylindrical shrinker ends together, including quadratic opening modes for which ordinary analytic graph arguments do not suffice. Finite Taylor jets and exponentially small end terms produce a localized gradient inequality. Increasing graph radii are chosen from preceding energy drops so that cutoff errors are summable; finite rescaled length gives convergence with its ambient position and axes fixed.

\subsection{Scope relative to this problem}
\textit{Uniqueness of surface tangent flow}: The theorem covers all first-time singularities of smooth compact embedded surfaces in R3, while the catalogue statement includes general mean-curvature-flow singularities, later times, dimensions and a further asymptotic-symmetry program.

This page records partial support. Complete resolution of the catalogue question is not established by this record.

\subsection{Formal verification}
\textit{Uniqueness of surface tangent flow}: The pinned catalogue lists no Lean formalization for this result.

This site has not run the Lean development. These links describe OpenAI’s selected formal statements and should be read with their stated scope.

\subsection{Status caveat}
OpenAI reports that an internal model produced the manuscripts. Their arguments have not been independently verified as proofs by this site. Any portion outside the documented formal scope may contain mistakes, and the catalogue label reports the announced mathematical claim.

\subsection{Sources}
\begin{itemize}
\item OpenAI, \emph{Uniqueness of tangent flows at the first singular time of embedded surface mean-curvature flow}, 2026-09-24; citation key \texttt{OAI:Tangent-flow-uniqueness-2026-09-24}. \href{https://github.com/openai/math/blob/adc7f1241b42e322a6451854ab7e4b4c146bf78a/preprints/Tangent-flow-uniqueness-2026-09-24/paper.pdf}{PDF (GitHub)}; \href{https://github.com/openai/math/raw/adc7f1241b42e322a6451854ab7e4b4c146bf78a/preprints/Tangent-flow-uniqueness-2026-09-24/paper.pdf}{PDF (direct)}; \href{https://github.com/openai/math/blob/adc7f1241b42e322a6451854ab7e4b4c146bf78a/preprints/Tangent-flow-uniqueness-2026-09-24/README.md}{Manuscript page and citation}; \href{https://github.com/openai/math/blob/main/preprints/Tangent-flow-uniqueness-2026-09-24/paper.pdf}{Latest version}.
\item \href{https://github.com/openai/math/blob/adc7f1241b42e322a6451854ab7e4b4c146bf78a/CONTENTS.md}{OpenAI catalogue entry}.
\end{itemize}
\end{document}
